\subsection{Absolutely Semisimple Algebras}

DEFINITION 2. --- Let K be a commutative field. We say that a K-algebra A is absolutely semisimple if the ring \(\mathrm{A}_{(\mathrm{L})}\) is semisimple for every extension L of K.

An absolutely semisimple algebra is semisimple. The K-algebra A is absolutely semisimple if and only if the A-module \(A_{s}\) is absolutely semisimple. By Proposition 1 of VIII, p.~229, we therefore obtain the following result: if

L is an extension of K, then the L-algebra \(A_{(L)}\) is absolutely semisimple if and only if the K-algebra A is absolutely semisimple.

THEOREM 1. --- Let K be a commutative field and A a K-algebra. The following properties are equivalent:

\begin{enumerate}
\def\labelenumi{(\roman{enumi})}
\item
  The K-algebra A is absolutely semisimple.
\item
  The algebra \(\mathrm{A}\) has finite degree over \(\mathrm{K}\) , and there exists an extension \(\mathrm{P}\) of \(\mathrm{K}\) that is a perfect field, for which the \(\mathrm{P}\) -algebra \(\mathrm{A}_{(\mathrm{P})}\) is semisimple.
\item
  The K-algebra A is semisimple and has finite degree over K, and its center is an étale K-algebra.
\item
  There exists a finite family \((n_{i},\mathrm{D}_{i})_{i\in\mathrm{I}}\) , where \(n_{i}\) is a strictly positive integer and \(D_{i}\) a K-algebra of finite degree that is a field, such that the center \(Z_{i}\) of \(D_{i}\) is a separable extension of K for every \(i\in I\) and A is isomorphic to the product of the matrix rings \(\mathbf{M}_{n_{i}}(\mathrm{D}_{i})\) .
\item
  There exist an extension \(\mathrm{L}\) of \(\mathrm{K}\) and a finite family of integers \((n_i)_{i\in \mathrm{I}}\) such that the L-algebra \(\mathrm{A}_{(\mathrm{L})}\) is isomorphic to the algebra \(\prod_{i\in \mathrm{I}}\mathbf{M}_{n_i}(\mathrm{L})\) .
\item
  There exist a Galois extension L of K of finite degree and a finite family of integers \((n_{i})\) such that \(\mathrm{A}_{(\mathrm{L})}\) is isomorphic to the product of the matrix algebras \(\mathbf{M}_{n_{i}}(\mathrm{L})\) .
\end{enumerate}

We first prove the implications (v) \(\Rightarrow\) (iv) \(\Rightarrow\) (iii) \(\Rightarrow\) (ii) \(\Rightarrow\) (i). Denote the center of A by Z.

If property (v) holds, then the L-algebra \(\mathrm{A}_{(\mathrm{L})}\) is semisimple and has finite degree over L, and its center (isomorphic to \(\mathrm{Z}_{(\mathrm{L})}\) , III, p.~41, Corollary) is isomorphic to \(L^{r}\) for some integer \(r \geqslant 0\) . By Corollary 2, a) of VIII, p.~222, the algebra A is semisimple and has finite degree over K. It is therefore isomorphic to a finite product of rings \(\prod_{i \in I} M_{n_{i}}(D_{i})\) with \(n_{i} \geqslant 1\) for every \(i \in I\) , where the field \(D_{i}\) is an algebra of finite degree over K. The center \(Z_{i}\) of \(D_{i}\) is a commutative field that is an extension of K, and Z is isomorphic to \(\prod_{i \in I} Z_{i}\) . Therefore, \(\mathrm{Z}_{(\mathrm{L})}\) is isomorphic, on the one hand, to \(\prod_{i \in I} (\mathrm{Z}_{i})_{(\mathrm{L})}\) and, on the other hand, to \(L^{r}\) . In other words, the algebra \(\prod_{i \in I} Z_{i}\) is étale over the field K (V, §6, No.~3, p.~28, Definition 1), and each of the extensions \(Z_{i}\) is separable over K (V, §6, No.~4, p.~31, Corollary, and V, §7, No.~1, p.~42). So (v) implies (iv).

If property (iv) holds, then it is clear that A is a semisimple algebra and has finite degree over K. Its center Z is isomorphic to the product \(\prod_{i\in I}Z_{i}\) of separable extensions of K of finite degree; it is therefore an étale algebra (loc. cit.). So (iv) implies (iii).

The implications (iii) \(\Rightarrow\) (ii) \(\Rightarrow\) (i) follow from Proposition 2 of VIII, p.~230 applied to the A-module \(\mathrm{A}_s\) .

Lemma 2. --- Let L be an algebraically closed field and D a field containing L in its center. If D is distinct from L, then there exists an extension \(L'\) of L such that the ring \(D \otimes_{L} L'\) is not left Artinian.

Let \(x\) be an element of \(\mathrm{D} - \mathrm{L}\) ; since \(\mathrm{L}\) is algebraically closed, the extension \(\mathrm{L}' = \mathrm{L}(x)\) of \(\mathrm{L}\) is not algebraic and \(x\) is transcendent over \(\mathrm{L}\) . The ring \(\mathrm{B} = \mathrm{L}' \otimes_{\mathrm{L}} \mathrm{L}'\) is then an integral domain by Proposition 5 of \(\mathrm{V}\) , §17, No.~4, p.~141.

The element \(y = x \otimes 1 - 1 \otimes x\) of B is not zero, but if \(\varphi\) is the homomorphism from B to \(L'\) that sends \(\xi \otimes \eta\) to \(\xi \eta\) , then we have \(\varphi(y) = 0\) , so y is not invertible in B. We view the ring \(C = D \otimes_{L} L'\) as a right module over its subring B; it is a free module because D is a right vector space over its subfield \(L'\) . Since y is a nonzero and noninvertible element of the integral domain B, right multiplication by y in C is a mapping \(R_y\) that is injective but not bijective. Now, \(R_y\) is an endomorphism of the left C-module \(C_s\) ; consequently (VIII, p.~28, Corollary 1), the ring C is not left Artinian.

Let us now prove that (i) implies (v). This is a consequence of the following lemma.

Lemma 3. --- Let A be an absolutely semisimple K-algebra and L an algebraically closed extension of K. Then the algebra \(\mathrm{A}_{(\mathrm{L})}\) is isomorphic to a product of finitely many matrix algebras over L.

The L-algebra \(\mathrm{A}_{(\mathrm{L})}\) is semisimple; it is therefore isomorphic to a product of finitely many algebras of the form \(\mathbf{M}_{n_{i}}(\mathrm{D}_{i})\) , where \(D_{i}\) is a field containing L in its center and \(n_{i}\) an integer \(\geqslant 1\) (VIII, p.~137, Remark 1).

Let \(L'\) be an extension of L. Since the K-algebra A is absolutely semisimple, the ring \(A_{(L')}\) is semisimple and therefore left Artinian. Now, the ring \(A_{(L')}\) is isomorphic to \(L' \otimes_L A_{(L)}\) , hence to \(\prod_{i \in I} M_{n_i}(L' \otimes_L D_i)\) ; by Proposition 5 of VIII, p.~7, each of the rings \(M_{n_i}(L' \otimes_L D_i)\) is therefore left Artinian.

Let \(n \geqslant 1\) be an integer and B a ring. Let \((\mathfrak{b}_r)_{r \geqslant 0}\) be a decreasing sequence of left ideals of B; denote by \(\mathfrak{c}_r\) the set of square matrices of order \(n\) with entries in \(\mathfrak{b}_r\) . Then \((\mathfrak{c}_r)_{r \geqslant 0}\) is a decreasing sequence of left ideals of \(\mathbf{M}_n(\mathrm{B})\) . In particular, if the ring \(\mathbf{M}_n(\mathrm{B})\) is left Artinian, then so is B.

By the above, for every \(i \in I\) and every extension \(L'\) of L, the ring \(D_{i} \otimes_{L} L'\) is left Artinian. By Lemma 2, we have \(D_{i} = L\) for every \(i \in I\) , which implies Lemma 3.

We will use the following lemma to prove the implication (i)⇒(vi).

Lemma 4. --- Let A and B be algebras over the field K that have finite generating sets, and let \(K'\) be an extension of K. If the \(K'\) -algebras \(\mathrm{A}_{(\mathrm{K}')}\) and \(\mathrm{B}_{(\mathrm{K}')}\) are isomorphic, then there exists a subextension L of \(K'\) , finitely generated over K, such that the L-algebras \(\mathrm{A}_{(\mathrm{L})}\) and \(\mathrm{B}_{(\mathrm{L})}\) are isomorphic.

Let \((e_{i})_{i\in\mathrm{I}}\) and \((f_{j})_{j\in\mathrm{J}}\) be finite generating sets for the algebras A and B, respectively. Let u be an isomorphism from \(\mathrm{A}_{(\mathrm{K}^{\prime})}\) to \(\mathrm{B}_{(\mathrm{K}^{\prime})}\) and v the inverse isomorphism; there exists a subextension L of \(K^{\prime}\) , finitely generated over K, such that we have \(u(1\otimes e_{i})\in\mathrm{B}_{(\mathrm{L})}\) for every \(i\in\mathrm{I}\) and \(v(1\otimes f_{j})\in\mathrm{A}_{(\mathrm{L})}\) for every \(j\in\mathrm{J}\) . Consequently, u sends \(\mathrm{A}_{(\mathrm{L})}\) into \(\mathrm{B}_{(\mathrm{L})}\) , and v sends \(\mathrm{B}_{(\mathrm{L})}\) into \(\mathrm{A}_{(\mathrm{L})}\) . The induced mappings \(u^{\prime}:\mathrm{A}_{(\mathrm{L})}\to\mathrm{B}_{(\mathrm{L})}\) and \(v^{\prime}:\mathrm{B}_{(\mathrm{L})}\to\mathrm{A}_{(\mathrm{L})}\) are ring homomorphisms; they are inverse bijections.

Let us complete the proof of the implication \((\mathrm{i})\Rightarrow(\mathrm{vi})\) . Let \(K'\) be a separable closure of K (V, §7, No.~8, p.~45). Then \(A_{(K')}\) is an absolutely semisimple algebra over \(K'\) . By the implication \((\mathrm{i})\Rightarrow(\mathrm{iv})\) , the \(K'\) -algebra \(A_{(K')}\) is isomorphic to a product \(\mathbf{M}_{n_{1}}(\mathrm{D}_{1})\times\cdots\times\mathbf{M}_{n_{r}}(\mathrm{D}_{r})\) , where \(D_{i}\) is a \(K'\) -algebra of finite degree that is a field whose center \(Z_{i}\) is a separable extension of \(K'\) . Since \(K'\) is separably closed, we have \(Z_{i}=K'\) . By Proposition 3 of VIII, p.~231, the field \(D_{i}\) is commutative. We therefore have \(D_{i}=K'\) . We denote the K-algebra \(\mathbf{M}_{n_{1}}(\mathrm{K})\times\cdots\times\mathbf{M}_{n_{r}}(\mathrm{K})\) by B. The \(K'\) -algebras \(A_{(K')}\) and \(B_{(K')}\) are isomorphic by the above. Every subextension of \(K'\) that is finitely generated over K is separable and has finite degree over K, hence is contained in a subextension L of \(K'\) that is Galois and of finite degree over K (V, §10, No.~1, p.~57, Proposition 2). The implication therefore follows from Lemma 4.

The implication (vi)⇒(v) is immediate.

COROLLARY 1. --- Let K be a commutative field, and let \(A_{1}\) and \(A_{2}\) be K-algebras. Suppose that \(A_{1}\) is absolutely semisimple.

\begin{enumerate}
\def\labelenumi{\alph{enumi})}
\item
  If \(\mathrm{A}_2\) is semisimple, then so is \(\mathrm{A}_1\otimes_{\mathrm{K}}\mathrm{A}_2\)
\item
  If \(\mathrm{A}_2\) is absolutely semisimple, then so is \(\mathrm{A}_1\otimes_{\mathrm{K}}\mathrm{A}_2\) .
\end{enumerate}

Denote the center of \(A_{1}\) by \(Z_{1}\) and that of \(A_{2}\) by \(Z_{2}\) . The center Z of \(A_{1} \otimes_{K} A_{2}\) is equal to \(Z_{1} \otimes_{K} Z_{2}\) by the corollary of III, §4, No.~4, p.~468. Suppose that \(A_{2}\) is semisimple; then \(Z_{2}\) is a reduced algebra (VIII, p.~137, Propositions 2 and 3). By Theorem 1, \(Z_{1}\) is an étale and therefore separable

K-algebra. By Proposition 5 of V, §15, No.~2, p.~120, the ring \(Z = Z_{1} \otimes_{K} Z_{2}\) is reduced; since \(A_{1}\) has finite degree over K (Theorem 1), it follows from Proposition 7 of VIII, p.~221 that the ring \(A_{1} \otimes_{K} A_{2}\) is semisimple.

Now suppose that \(A_{2}\) is absolutely semisimple. Let L be an extension of K. Then the algebra \(\mathrm{A}_{1(\mathrm{L})}\) is absolutely semisimple, and the algebra \(\mathrm{A}_{2(\mathrm{L})}\) is semisimple. Therefore, by a), the algebra \(\mathrm{A}_{1} \otimes_{\mathrm{K}} \mathrm{A}_{2(\mathrm{L})}\) is semisimple.

COROLLARY 2. --- Let K be a separably closed field, and let A be an absolutely semisimple K-algebra. Then there exist an integer \(r \geqslant 0\) and strictly positive integers \(n_{1}, \ldots, n_{r}\) such that the algebra A is isomorphic to the algebra \(\prod_{i=1}^{r} \mathbf{M}_{n_{i}}(\mathrm{K})\) .

By Theorem 1, A is isomorphic to an algebra of the form \(\prod_{i=1}^{r} \mathbf{M}_{n_i}(\mathrm{D}_i)\) for an integer \(r \geqslant 0\) , integers \(n_1, \ldots, n_r\) , and K-algebras of finite degree \(\mathrm{D}_1, \ldots, \mathrm{D}_r\) that are fields and whose centers are separable extensions of K and therefore equal to K. By Proposition 3 of VIII, p.~231, we have \(\mathrm{D}_i = \mathrm{K}\) for \(i \in [1, r]\) .

Example. --- A commutative K-algebra is absolutely semisimple if and only if it is étale: this follows from the definition (V, §6, No.~3, p.~28, Definition 1) and the equivalence of properties (i) and (v) of Theorem 1.

